\section*{अध्याय \ref{ch:b2:corrosion} --- संक्षारण और संरक्षण}

\begin{solution}{exo:b2:corrosion:1}
$j = \qty{0.020}{A/m^2}$: $de/dt = 0.020 \times 65.38/(2 \times 96\,485 \times 7.13
\times 10^6) = \qty{9.5e-13}{m/s}$, अर्थात् प्रति वर्ष \qty{0.030}{mm}।
\end{solution}

\begin{solution}{exo:b2:corrosion:2}
ताँबे पर इस्पात के बोल्ट: लोहा, कम उत्कृष्ट, ऐनोड है, और बड़े ताँबा कैथोड पर
छोटे बोल्ट तेज़ी से संक्षारित होते हैं। इस्पात में ज़िंक की कीलें: ज़िंक संक्षारित होता है
और अपने चारों ओर के इस्पात की रक्षा करता है।
\end{solution}

\begin{solution}{exo:b2:corrosion:3}
प्लेटों के बीच के संकरे अंतराल में और बोल्ट के सिर के नीचे, जहाँ जल
धीरे-धीरे नवीकृत होता है और उसकी ऑक्सीजन जल्दी समाप्त हो जाती है: \omterm{def:b2:corrosion:differential}{विभेदी वातन} से
ये बंद क्षेत्र ऐनोड बन जाते हैं, खुली सतहें कैथोड।
\end{solution}

\begin{solution}{exo:b2:corrosion:4}
टिन का डिब्बा: लोहे से अधिक उत्कृष्ट टिन उघड़े इस्पात को एक बड़े कैथोड का
ऐनोड बना देता है। यशदलेपित बाल्टी पर इस्पात के बजाय ज़िंक
संक्षारित होता है।
\end{solution}

\begin{solution}{exo:b2:corrosion:5}
विलोडन \omterm{def:b2:current-potential-curves:limiting-current}{विसरण परत} को पतला करता है: ऑक्सीजन पठार ऊपर उठता है, लोहे का ऐनोडी वक्र
उससे ऊँचाई पर मिलता है, अतः \omterm{def:b2:corrosion:uniform}{संक्षारण विभव} और \omterm{def:b2:corrosion:uniform}{संक्षारण
धारा} दोनों बढ़ते हैं। ऑक्सीजन के बिना उदासीन जल में बची एकमात्र कैथोडी अभिक्रिया
जल का मंद अपचयन है: \omterm{def:b2:corrosion:uniform}{संक्षारण विभव} गिरता है और
\omterm{def:b2:corrosion:uniform}{संक्षारण धारा} बहुत छोटी हो जाती है।
\end{solution}

\begin{solution}{exo:b2:corrosion:6}
$t = 0.90 \times 5000 \times 2 \times 96\,485/(65.38 \times 0.10) = \qty{1.33e8}{s}$,
लगभग 4.2~वर्ष।
\end{solution}

\begin{solution}{exo:b2:corrosion:7}
प्रति किलोग्राम आवेश, $nF/M$: ज़िंक \qty{820}{A.h/kg}, मैग्नीशियम \qty{2.21e3}{A.h/kg},
ऐलुमिनियम \qty{2.98e3}{A.h/kg}। मानक विभव: \ce{Zn^2+}/\ce{Zn} $-0.76$,
\ce{Mg^2+}/\ce{Mg} $-2.36$, \ce{Al^3+}/\ce{Al} \qty{-1.68}{V}। सूखी, प्रतिरोधी मिट्टी में
धारा को बड़े प्रतिरोध से होकर धकेलना पड़ता है: मैग्नीशियम, लोहे से कहीं
नीचे होने के कारण, सबसे बड़ी चालक वोल्टता देता है।
\end{solution}

\begin{solution}{exo:b2:corrosion:8}
$U = 2.0 \times 4.0 + 1.5 = \qty{9.5}{V}$; शक्ति \qty{19}{W}; प्रति वर्ष $19 \times
8766 = \qty{1.7e5}{W.h}$, \qty{1.7e2}{kW.h}।
\end{solution}

\begin{solution}{exo:b2:corrosion:9}
जल-रेखा पर धातु ऑक्सीजन-समृद्ध जल से गीली रहती है, एक कैथोड; ठीक
नीचे जल कम वातित है और वह क्षेत्र जल-रेखा के साथ बने सेल का ऐनोड बन जाता है:
इस्पात वहाँ एक पट्टी में संक्षारित होता है।
\end{solution}

\begin{solution}{exo:b2:corrosion:10}
वातित जल में स्टेनलेस इस्पात का \omterm{def:b2:corrosion:uniform}{संक्षारण विभव} उसके
निष्क्रिय पठार पर पड़ता है: नगण्य धारा, परत स्वयं भरती रहती है। क्लोराइड आयन
कमज़ोर बिंदुओं पर परत को तोड़ते हैं; वहाँ नंगी धातु एक विशाल निष्क्रिय कैथोड के
सामने छोटा ऐनोड है, इसलिए उसका धारा घनत्व बहुत बड़ा होता है। गर्त के भीतर
धातु आयन जल-अपघटित होकर विलयन को अम्लीय बनाते हैं, आवेश संतुलित करने के लिए क्लोराइड भीतर
खिंचता है, और परत फिर नहीं बन पाती: गर्त गहरा होता जाता है।
\end{solution}

\begin{solution}{exo:b2:corrosion:11}
ऑक्सीजन पूरे \qty{22}{cm^2} पर अपचित होती है: \omterm{def:b2:current-potential-curves:ie-curve}{कैथोडी धारा} $22 \times 5 =
\qty{110}{\micro A}$, जो पूरी की पूरी जोड़ के निकट के \qty{2}{cm^2} इस्पात के ऑक्सीकरण
से आती है, $j = \qty{55}{\micro A/cm^2}$। हानि \qty{10}{\micro A/cm^2} पर लोहे की हानि
की $5.5$ गुनी है, प्रति वर्ष लगभग \qty{0.64}{mm}।
\end{solution}

\begin{solution}{exo:b2:corrosion:12}
लोहा \ce{Fe^2+}/\ce{Fe} रेखा से नीचे प्रतिरक्षित है, $-0.41 + 0.0296\log 10^{-6} =
\qty{-0.59}{V}$: लगभग \qty{-0.6}{V} से नीचे रखने पर इस्पात अब नहीं घुलता। बहुत
नीचे, इस्पात पर जल बढ़ती दर से अपचित होता है (\ce{H+}/\ce{H2} रेखा pH 7 पर
\qty{-0.41}{V} पर है और इस्पात पर \omterm{def:b2:current-potential-curves:overpotential}{अधिविभव} वोल्ट के कुछ दसवें भाग है):
धारा व्यर्थ जाती है और हाइड्रोजन लेप को उखाड़ देती है।
\end{solution}

\begin{solution}{pb:b2:corrosion:1}
\textbf{1.} \ce{Fe -> Fe^2+ + 2e-}; \ce{O2 + 2H2O + 4e- -> 4OH-}।
\textbf{2.} ऑक्सीजन का अपचयन, मंद और अपने विसरण पठार तक पहुँचा हुआ,
\omterm{def:b2:current-potential-curves:ie-curve}{कैथोडी धारा} को सीमित करता है; लोहे का ऐनोडी वक्र उससे उसी पठार पर मिलता है।
\textbf{3.} $j = \qty{0.10}{A/m^2}$: $0.10 \times 3.156 \times 10^7/(2 \times 96\,485)
\times 55.845 = \qty{913}{g}$ प्रति वर्ग मीटर और प्रति वर्ष।
\textbf{4.} $913/7.87 \times 10^6 = \qty{1.16e-4}{m}$, प्रति वर्ष \qty{0.12}{mm}।
\textbf{5.} $4/0.116 = 34$~वर्ष।
\textbf{6.} संक्षारण एकसमान नहीं है: दरारें, मिट्टियों के बीच वातन के
अंतर, और लेप के दोष \omterm{def:b2:current-potential-curves:ie-curve}{ऐनोडी धारा} को छोटे क्षेत्रों पर केंद्रित करते हैं,
जहाँ गर्त बहुत पहले दीवार में छेद कर देते हैं।
\textbf{7.} हर युग्म के लिए $E^\circ = \Delta_f G^\circ/(nF)$: \ce{Fe} $-0.41$, \ce{Zn} $-0.76$,
\ce{Mg} $-2.36$, \ce{Al} \qty{-1.68}{V}।
\textbf{8.} ज़िंक, मैग्नीशियम और ऐलुमिनियम, सभी लोहे से नीचे।
\textbf{9.} ऐलुमिनियम निष्क्रिय हो जाता है: उसकी ऑक्साइड परत उसका विलयन रोक देती है और
तब वह कम धारा देता है, जब तक परत को रोकने के लिए उसे मिश्रधातु न बनाया जाए (जो
समुद्री जल में काम करता है, मिट्टियों में नहीं)।
\textbf{10.} धारा को मिट्टी का प्रतिरोध पार करना पड़ता है: ऐनोड और संरक्षित
इस्पात के बीच चालक वोल्टता मैग्नीशियम के लिए लगभग \qty{1.9}{V} है, जबकि
ज़िंक के लिए \qty{0.35}{V} (मानक विभवों से)।
\textbf{11.} $3.156 \times 10^7 \times 24.305/(2 \times 96\,485 \times 0.50) =
\qty{7.95}{kg}$ प्रति ऐम्पियर और प्रति वर्ष।
\textbf{12.} $\pi \times 0.30 \times 10\,000 = \qty{9.42e3}{m^2}$।
\textbf{13.} \qty{94.2}{m^2}।
\textbf{14.} $0.020 \times 94.2 = \qty{1.88}{A}$।
\textbf{15.} $1.885 \times 20 \times 3.156 \times 10^7 = \qty{1.19e9}{C}$।
\textbf{16.} $1.19 \times 10^9 \times 24.305/(2 \times 96\,485 \times 0.50) =
\qty{3.00e5}{g}$, \qty{300}{kg} (या $1.885 \times 20 \times 7.95$)।
\textbf{17.} $300/15 = 20$ ऐनोड।
\textbf{18.} मिट्टी का प्रतिरोध सीमित करता है कि एक अकेला ऐनोड अपनी धारा को पाइप के अनुदिश
कितनी दूर तक धकेल सकता है: वितरित ऐनोड उसकी एकसमान रक्षा करते हैं।
\textbf{19.} एक दिष्टकारी भूमि-संस्तर में गड़े एक अक्रिय ऐनोड से, मिट्टी से होकर,
अपने ऋणात्मक ध्रुव से जुड़े पाइप में धारा चलाता है।
\textbf{20.} $1.885 \times 5.0 + 1.5 = \qty{10.9}{V}$।
\textbf{21.} $10.9 \times 1.885 = \qty{20.6}{W}$; $20.6 \times 8766 = \qty{1.8e5}{W.h}$,
प्रति वर्ष \qty{180}{kW.h}।
\textbf{22.} पाइप पर जल अपचित होता है: हाइड्रोजन बनती है, लेप को उखाड़ती है और
इस्पात को भंगुर बना सकती है, और धारा व्यर्थ जाती है।
\textbf{23.} $E < -0.41 + 0.0296\log(10^{-6}) = \qty{-0.59}{V}$।
\textbf{24.} बीस वर्षों के संरक्षण के लिए $\boldsymbol{\approx \qty{3.0e2}{kg}}$
मैग्नीशियम ऐनोड चाहिए।
\end{solution}
